% Generated draft: requires mathematical and source review.
\documentclass[11pt]{article}
\usepackage[margin=25mm]{geometry}
\usepackage{fontspec}
\setmainfont{Latin Modern Roman}
\usepackage{amsmath,amssymb,mathtools,unicode-math}
\setmathfont{Latin Modern Math}
\usepackage{enumitem,needspace,xurl,hyperref}
\setlength{\parindent}{0pt}
\setlength{\parskip}{5pt}
\setlength{\emergencystretch}{3em}
\newcommand{\N}{\mathbb{N}}
\newcommand{\Z}{\mathbb{Z}}
\newcommand{\Q}{\mathbb{Q}}
\newcommand{\R}{\mathbb{R}}
\newcommand{\C}{\mathbb{C}}
\newcommand{\F}{\mathbb{F}}
\newcommand{\A}{\mathbb{A}}
\newcommand{\PP}{\mathbb{P}}
\newcommand{\E}{\mathbb{E}}
\newcommand{\eps}{\varepsilon}
\newcommand{\dd}{\,\mathrm{d}}
\begin{document}
\section*{UnsolvedMath ID: 20002094}
\textbf{Source problem code:} catalogue entry
\subsection{Definitions and mathematical statement}
Question 1.36 asks when a conformally symplectic structure is global, stating that this is equivalent to exactness of a one-form denoted \(\omega\), and asking for other conditions. The paired aid interprets this as the Lee form. Let \(M^{2n}\) be connected, \(2n\geq4\), and let \(\Omega\) be nondegenerate with closed Lee form \(\theta\), using the source convention
\[
 d\Omega=-\theta\wedge\Omega,\qquad
 d_\theta=d+\theta\wedge.
\]
Then \(d_\theta\Omega=0\). Global conformality means there is a smooth \(f\) with \(e^f\Omega\) symplectic. The reported equivalent tests are
\[
 \theta=df\text{ globally},\qquad [\theta]=0\in H^1(M;\mathbb R),\qquad
 \rho_\theta(\gamma)=\int_\gamma\theta=0\quad\forall\gamma\in\pi_1(M).
\]
These use the specified sign and distinguish the Lee form from the nondegenerate two-form.

The aid gives two additional closed-manifold tests. Let \(\pi:\widetilde M\to M\) be the minimal Lee cover, corresponding to \(\ker\rho_\theta\), choose \(df=\pi^*\theta\), and put \(\widetilde\omega=e^f\pi^*\Omega\). For closed \(M\), globality is equivalent to finite total symplectic volume of \((\widetilde M,\widetilde\omega)\). Nonzero Lee periods produce deck homotheties, whose iterates force infinite total volume. Also no finite cover can globalize a strict lcs form, since finite-cover pullback is injective on real first cohomology.

On a closed manifold, twisted exactness \(\Omega=d_\theta\eta\) obstructs globality: if \(\theta=df\), then \(e^f\Omega=d(e^f\eta)\) would be an exact symplectic form on a closed manifold. This implication is not asserted on a noncompact base. Nonvanishing of \([\Omega]\in H^2_\theta(M)\), however, is not sufficient for globality. The source historically mentions examples of Banyaga with nonzero Lichnerowicz class, without a full reference.

For a mapping torus
\[
 T_\varphi=(N\times[0,1])/((x,1)\sim(\varphi(x),0)),
\]
with compact fiber \(N\) and Lee class \(b[dt]\), choose the representative \(\theta=b\,dt\). The aid reports a twisted Wang/mapping-cone description governed by
\[
 e^b\varphi^*-I.
\]
In real cohomology its degree-two part has the short exact form
\[
 0\longrightarrow\operatorname{coker}(e^b\varphi^*-I|_{H^1(N)})
 \longrightarrow H^2_{b\,dt}(T_\varphi)
 \longrightarrow\ker(e^b\varphi^*-I|_{H^2(N)})
 \longrightarrow0.
\]
Consequently a nonzero \([\Omega]\) requires \(e^{-b}\) in the monodromy spectrum on \(H^1(N;\mathbb R)\) or \(H^2(N;\mathbb R)\). Reversing the gluing or twisted-differential convention changes the dictionary, so the signs are not dropped.

The aid reports an explicit hyperbolic \(T^3\) mapping-torus example realizing this resonance with a nonzero twisted lcs class and nonzero Lee class. Its matrix and form are not supplied here, so it is retained as an attributed counterexample to sufficiency of twisted-class nonvanishing, not reconstructed. These are globality and resonance tests, not a classification of all conformally symplectic structures.
\subsection{Short English statement}
The supplied aid interprets globality as vanishing Lee class. On closed manifolds it adds a minimal-Lee-cover volume test and a twisted-exactness obstruction. Its mapping-torus formula imposes a sign-sensitive monodromy resonance for nonzero Lichnerowicz class, while a reported strict example shows that such nonvanishing alone does not imply globality.
\subsection{Sources}
\begin{itemize}
\item Local source record: catalogue entry, catalogue entry; canonical \texttt{aim-geometry-notes.json}, \texttt{notes[431]}, Question 1.36. Title: ``Globality, Lee-cover volume, and mapping-torus resonance for lcs forms.''
\item AIM workshop list \emph{Holomorphic curves in contact geometry}, supplied reference \url{https://aimath.org/WWN/contactgeom2/contactgeom2.pdf}. The supplied verification locates the text in sixth-day notes; the embedded page number is an extraction artifact. This draft does not access a PDF or HTML source.
\item Original attribution to Banyaga for nonzero Lichnerowicz-class examples, without a full citation. Attempt 1 supplies the Lee-cover/finite-cover tests, twisted Wang sequence and hyperbolic torus mapping-torus example; the latter's explicit data are not reproduced.
\item Registry dated 14 August 2026: ``partially\_solved,'' ``PARTIAL-PROGRESS,'' statement status ``exact''; attempt partial, rigor ``proved,'' novelty confidence low. The machine-generated, non-peer-reviewed aid requires independent verification before citation.
\end{itemize}
\end{document}
